
\begin{proposition}[The canonical normal-crossings formula]\label{cdp:st:A.1.2}
For a reduced normal-crossings divisor, the map in (A.3) is an isomorphism. In a chart with $F=x_1\cdots x_r$, it is explicitly
\[
 \bigoplus_{i=1}^r R/(x_i)\longrightarrow K,
 \qquad (a_i)_i\longmapsto
 \sum_{i=1}^r a_i\prod_{j\ne i}x_j\,dx_i.
\tag{A.6}\label{cdp:eq:A.6}
\]
The tensor factor $N^*$ is trivialized here by $[F]$. At a double crossing $T\simeq R/(x,y)$; at a triple crossing,
\[
 T\simeq R^3/\bigl(R(1,1,1)+R(x,0,0)+R(0,y,0)+R(0,0,z)\bigr).
\tag{A.7}\label{cdp:eq:A.7}
\]
\end{proposition}
\begin{proof}
On the normalization branch $x_i=0$, an ambient form maps to zero in the cotangent bundle of that branch precisely when all coefficients except that of $dx_i$ vanish. Thus a coefficient of $dx_i$ must vanish on every branch $x_j=0$ for $j\ne i$. In the convergent power-series ring it is divisible by $\prod_{j\ne i}x_j$. Its remaining factor is defined modulo $x_i$, which gives exactly \eqref{cdp:eq:A.6}. The coefficients of different $dx_i$ show injectivity. Coefficients of the coordinates not among $x_1,\ldots,x_r$ vanish on every branch and hence vanish in $R$. This proves surjectivity in every dimension.

For $r=1$ the map is $a\mapsto a\,dx$. For $r=2$ it is $(a,b)\mapsto ay\,dx+bx\,dy$, and quotienting by $dF=y\,dx+x\,dy$ leaves the generator $y\,dx=-x\,dy$, annihilated by $(x,y)$. For $r=3$ the branch generators are $yz\,dx,xz\,dy,xy\,dz$. Their sum is $dF$ and their respective annihilators are $(x),(y),(z)$, proving \eqref{cdp:eq:A.7}. Multiplication of $F$ by a unit and changes of coordinates were checked in \ref{cdp:st:A.2.1}. Formula \eqref{cdp:eq:A.6} is the restriction of the canonical fractional-ideal map, so it glues without a chosen branch ordering, equation, or coordinate system.
\end{proof}


